% META: {"id": "TNSI-ECRIT-S6-EX1", "chapitre": "TNSI-STRUCTURES-DONNEES", "type_objet": "banque_ecrite", "sujet": "TNSI-ECRIT-S6", "rang": 1, "capacites": ["T-STRUCT-05B", "T-STRUCT-05D"], "capacites_codes": ["C12", "C14"], "duree_min": 70, "autorite": "MENE2516123N", "mode_creation": "ex_nihilo", "sources_inspiration": [], "notice": "ORIGINAL_NEXUS_TRAINING_MATERIAL_NOT_OFFICIAL_EXAM_CONTENT", "status": "generated"}
% BEGIN-VERIFY
% matrice = [
%     [0, 1, 0, 1, 0],
%     [0, 0, 1, 0, 0],
%     [0, 0, 0, 0, 1],
%     [0, 0, 1, 0, 0],
%     [0, 0, 0, 0, 0],
% ]
% # Q1 : aretes et degres
% def aretes(m):
%     return [(i, j) for i, ligne in enumerate(m)
%             for j, v in enumerate(ligne) if v == 1]
% assert aretes(matrice) == [(0,1),(0,3),(1,2),(2,4),(3,2)]
% assert len(aretes(matrice)) == 5
% def degre_sortant(m, i):
%     return sum(m[i])
% def degre_entrant(m, j):
%     return sum(ligne[j] for ligne in m)
% assert [degre_sortant(matrice, i) for i in range(5)] == [2, 1, 1, 1, 0]
% assert [degre_entrant(matrice, j) for j in range(5)] == [0, 1, 2, 1, 1]
% assert sum(degre_sortant(matrice, i) for i in range(5)) == len(aretes(matrice))
% # Q2 : le graphe n'est pas symetrique, donc oriente
% assert any(matrice[i][j] != matrice[j][i]
%            for i in range(5) for j in range(5))
% assert matrice[0][1] == 1 and matrice[1][0] == 0
% # Q3 : passage aux listes de successeurs, et retour
% def matrice_vers_listes(m):
%     return {i: [j for j, v in enumerate(ligne) if v == 1]
%             for i, ligne in enumerate(m)}
% def listes_vers_matrice(listes):
%     n = len(listes)
%     m = [[0] * n for _ in range(n)]
%     for i, successeurs in listes.items():
%         for j in successeurs:
%             m[i][j] = 1
%     return m
% listes = matrice_vers_listes(matrice)
% assert listes == {0: [1, 3], 1: [2], 2: [4], 3: [2], 4: []}
% assert listes_vers_matrice(listes) == matrice
% # Q4 : cout de la question « i pointe-t-il vers j ? »
% cases = {"n": 0}
% def arete_matrice(m, i, j):
%     cases["n"] += 1
%     return m[i][j] == 1
% def arete_listes(listes, i, j):
%     for successeur in listes[i]:
%         cases["n"] += 1
%         if successeur == j:
%             return True
%     return False
% cases["n"] = 0; assert arete_matrice(matrice, 0, 3) is True and cases["n"] == 1
% cases["n"] = 0; assert arete_listes(listes, 0, 3) is True and cases["n"] == 2
% cases["n"] = 0; assert arete_matrice(matrice, 4, 0) is False and cases["n"] == 1
% cases["n"] = 0; assert arete_listes(listes, 4, 0) is False and cases["n"] == 0
% # Q5 : occupation memoire, graphe creux contre graphe dense
% def cases_matrice(n):
%     return n * n
% def cases_listes(listes):
%     return sum(len(s) for s in listes.values())
% assert cases_matrice(5) == 25
% assert cases_listes(listes) == 5
% # un graphe complet inverse le rapport
% complet = {i: [j for j in range(5) if j != i] for i in range(5)}
% assert cases_listes(complet) == 20 < cases_matrice(5)
% # et la matrice ne depend pas du nombre d'aretes
% vide = {i: [] for i in range(5)}
% assert cases_listes(vide) == 0
% assert cases_matrice(5) == 25
% END-VERIFY
\section*{Sujet S6 --- Exercice 1 : deux représentations d'un même graphe}

\textit{Durée conseillée : 70 minutes. Cet exercice est indépendant des deux
autres du sujet.}

\medskip
\textit{Contenu original Nexus --- ce n'est pas un sujet officiel d'examen.}

\subsection*{Énoncé}

Un graphe orienté à cinq sommets, numérotés de $0$ à $4$, est donné par sa
matrice d'adjacence. La case ligne $i$, colonne $j$ vaut $1$ s'il existe une
flèche $i \rightarrow j$.

\begin{python}
matrice = [
    [0, 1, 0, 1, 0],
    [0, 0, 1, 0, 0],
    [0, 0, 0, 0, 1],
    [0, 0, 1, 0, 0],
    [0, 0, 0, 0, 0],
]
\end{python}

\begin{enumerate}
  \item Lister toutes les flèches sous la forme $i \rightarrow j$. Donner, pour
    chaque sommet, son degré sortant et son degré entrant. Vérifier une relation
    entre la somme des degrés sortants et le nombre de flèches.

  \item Justifier, en citant deux cases précises, que ce graphe est
    \textbf{orienté} et non pas non orienté.

  \item Écrire \lstinline{matrice_vers_listes} et \lstinline{listes_vers_matrice}.
    Donner les listes de successeurs obtenues, et vérifier que la conversion
    inverse redonne la matrice de départ.

  \item Pour répondre à « le sommet $i$ pointe-t-il vers $j$ ? », on compte le
    nombre de cases examinées. Donner ce nombre, dans chaque représentation,
    pour $(0, 3)$ puis pour $(4, 0)$. Commenter le cas $(4, 0)$.

  \item Comparer l'occupation mémoire des deux représentations sur ce graphe,
    puis sur un graphe complet à cinq sommets, puis sur un graphe sans aucune
    flèche. En déduire dans quel cas chaque représentation est préférable.
\end{enumerate}

\subsection*{Corrigé}

\textbf{Question 1.} \emph{(4 points)}

Flèches : $0 \rightarrow 1$, $0 \rightarrow 3$, $1 \rightarrow 2$,
$2 \rightarrow 4$, $3 \rightarrow 2$. Il y en a $5$.

\begin{center}
\begin{tabular}{|c|c|c|c|c|c|}
\hline
\textbf{Sommet} & $0$ & $1$ & $2$ & $3$ & $4$ \\
\hline
\textbf{Degré sortant} & $2$ & $1$ & $1$ & $1$ & $0$ \\
\hline
\textbf{Degré entrant} & $0$ & $1$ & $2$ & $1$ & $1$ \\
\hline
\end{tabular}
\end{center}

La somme des degrés sortants vaut $5$, soit le nombre de flèches : chaque
flèche sort d'exactement un sommet. La somme des degrés entrants vaut aussi
$5$, pour la même raison.

\medskip
\textbf{Question 2.} \emph{(2 points)}

\lstinline{matrice[0][1]} vaut $1$ et \lstinline{matrice[1][0]} vaut $0$ : la
matrice n'est pas symétrique. Le sommet $0$ pointe vers $1$ sans que $1$ pointe
vers $0$.

\medskip
\textbf{Question 3.} \emph{(5 points)}

\begin{python}
def matrice_vers_listes(m):
    return {i: [j for j, valeur in enumerate(ligne) if valeur == 1]
            for i, ligne in enumerate(m)}

def listes_vers_matrice(listes):
    n = len(listes)
    m = [[0] * n for _ in range(n)]
    for i, successeurs in listes.items():
        for j in successeurs:
            m[i][j] = 1
    return m
\end{python}

Listes : \lstinline!{0: [1, 3], 1: [2], 2: [4], 3: [2], 4: []}!. La conversion
inverse redonne exactement la matrice de départ.

\medskip
\textbf{Question 4.} \emph{(5 points)}

\begin{center}
\begin{tabular}{|c|c|c|}
\hline
\textbf{Question} & \textbf{Matrice} & \textbf{Listes} \\
\hline
$(0, 3)$ & $1$ case & $2$ cases \\
\hline
$(4, 0)$ & $1$ case & $0$ case \\
\hline
\end{tabular}
\end{center}

La matrice répond toujours en une seule lecture, quelle que soit la question.
Les listes dépendent du degré sortant : pour $(4, 0)$, le sommet $4$ n'a aucun
successeur, la boucle ne s'exécute pas et la réponse est immédiate. Les listes
sont donc parfois plus rapides, jamais garanties de l'être.

\medskip
\textbf{Question 5.} \emph{(4 points)}

\begin{center}
\begin{tabular}{|l|c|c|}
\hline
\textbf{Graphe} & \textbf{Matrice} & \textbf{Listes} \\
\hline
celui de l'énoncé ($5$ flèches) & $25$ & $5$ \\
\hline
complet ($20$ flèches) & $25$ & $20$ \\
\hline
sans flèche & $25$ & $0$ \\
\hline
\end{tabular}
\end{center}

La matrice occupe $n^2$ cases \textbf{quel que soit} le nombre de flèches ; les
listes en occupent autant qu'il y a de flèches. Pour un graphe
\textbf{creux} — peu de flèches devant $n^2$ — les listes sont bien plus
économiques. Pour un graphe \textbf{dense}, l'écart se referme et la matrice
reprend l'avantage par la rapidité de son accès.

\subsection*{Critique}

\begin{itemize}
  \item \textbf{Lire la matrice en colonnes.} Confondre degré entrant et degré
    sortant est l'erreur la plus fréquente. La vérification par la somme, en
    question 1, permet au candidat de s'en apercevoir seul.
  \item \textbf{Croire les listes toujours meilleures.} La question 4 montre le
    contraire sur $(0, 3)$ : deux lectures au lieu d'une. Le cas $(4, 0)$ est là
    pour empêcher la conclusion inverse.
  \item \textbf{Oublier le sommet isolé dans la conversion.} Le sommet $4$ doit
    figurer dans les listes avec une liste \textbf{vide}. L'omettre casse la
    conversion inverse, et la vérification de la question 3 le révèle.
  \item \textbf{Compter la mémoire en « nombre de listes ».} Ce qui compte est
    le nombre total de successeurs stockés, pas le nombre d'entrées du
    dictionnaire — lequel vaut toujours $n$.
\end{itemize}

\subsection*{Portée pédagogique}

L'exercice fait constater qu'une même information admet deux représentations
exactes, et que le choix entre elles ne se juge pas dans l'absolu : il dépend
de la densité du graphe et de la question qu'on lui pose.

La conversion aller-retour de la question 3 fournit au candidat sa propre
vérification : si la matrice reconstruite diffère, l'une des deux fonctions est
fausse, et il le voit sans corrigé. Les questions 4 et 5 séparent deux critères
qu'on résume souvent en un seul mot « efficace » — le temps d'une requête et
la place occupée — et montrent qu'ils ne désignent pas le même gagnant.
